\section{Discussion: each axis bites at a different stage}\label{sec:discussion}

Table~\ref{tab:staged} re-estimates every outcome on one harmonized
sample --- one row per debut owner, scored on that owner's first filing
--- under a single specification, so the coefficients can be read down a
column. Magnitudes differ from the section-specific tables above.

\begin{table}[h]\centering\footnotesize
\caption{Every outcome in the funnel, conditioned on its prerequisite, under one specification. Unit is the debut owner, scored on that owner's first filing. Each row is a linear probability model with class$\times$debut-year fixed effects and heteroskedasticity-robust errors; coefficients are percentage points per standard deviation. \emph{Atypicality} is the average of the two KL levels; \emph{lead} is the signed difference; the lead-magnitude column $|L|$ comes from a second specification that adds it. Cohort windows differ because each maintenance deadline needs elapsed time and Form~D is electronic only from 2009, so the rows are not a nested decomposition. The five-year proof row is survival of a dated cancellation for non-use; the year-ten row is whether a registration that passed the proof was renewed rather than cancelled or expired, restricted to 2002--2013 cohorts, since a year-six and a year-ten cancellation share a terminal status and are separable only by date.}
\label{tab:staged}
\begin{tabular}{lrrrrr}
\toprule
 & & & \multicolumn{2}{c}{Unsigned} & Signed \\
\cmidrule(lr){4-5}\cmidrule(lr){6-6}
Outcome \emph{(population)} & $n$ & Base & Atyp. & $|L|$ & Lead \\
\midrule
Reached registration \emph{(filed)} & 3,161,552 & 56.11\% & $-0.982$ (0.032) & $-0.497$ (0.047) & $+0.475$ (0.028) \\
Passed the five-year proof \emph{(registered)} & 1,260,392 & 41.82\% & $+1.311$ (0.053) & $-0.527$ (0.083) & $-0.968$ (0.051) \\
Renewed at year ten \emph{(passed the proof)} & 297,944 & 55.92\% & $+1.383$ (0.104) & $-0.221$ (0.156) & $-0.428$ (0.096) \\
Raised a Reg D round \emph{(filed)} & 1,535,004 & 2.38\% & $-0.125$ (0.015) & $+0.345$ (0.027) & $-0.393$ (0.017) \\
Raised a Reg D round \emph{(registered)} & 916,448 & 2.45\% & $-0.169$ (0.020) & $+0.329$ (0.035) & $-0.324$ (0.022) \\
Listed after the round \emph{(funded)} & 36,475 & 3.51\% & $+0.866$ (0.138) & $+0.734$ (0.195) & $-0.534$ (0.129) \\
Became SEC-reporting \emph{(registered)} & 1,773,938 & 0.56\% & $+0.017$ (0.007) & $+0.023$ (0.012) & $-0.018$ (0.007) \\
\bottomrule
\end{tabular}
\\[0.3em]\footnotesize Standard errors in parentheses, in place of significance stars; see the note to Table~\ref{tab:gate_lpm}.
\end{table}

Surprising language costs at registration and pays afterwards: atypical
debut filings complete registration less often ($-1.35$pp per standard
deviation), then survive the five-year proof of continued use more often
($+1.35$pp) and reach public markets more often ($+0.04$pp on a 0.56\%
base --- the one specification in which a positive atypicality
coefficient on listing survives, and it is a continuous term through a
non-monotone profile). Having been early runs the other way: leading
applications complete registration \emph{more} often ($+0.87$pp) and then
pass the five-year proof \emph{less} often ($-0.49$pp on passing). Both
survival rows point the same way at the year-ten renewal, four years
later and on a sixth of the sample: atypicality $+1.30$pp and lead
$-0.63$pp among marks that had already cleared the five-year proof.
Whatever the five-year proof selects on, the renewal selects on it again
rather than reversing it.

In the financing rows atypicality does nothing at one stage and a great
deal at the next. Whether a debut owner ever raises a Regulation~D round
is unrelated to how atypical its filing was ($-0.034$pp per standard
deviation against a 2.45\% base, $t = -1.7$); among owners who did raise
one, the same measure predicts \emph{reaching} the public market
($+0.61$pp on a 3.51\% base, $t = 4.4$). Atypical language is invisible
at the financing stage and informative afterwards, which is what an
investor screen that does not read the text would look like --- the text
being free and public at the moment of filing. The comparison is among
survivors, so it cannot separate that reading from selection on
unobservables. Three limits bound these rows: Regulation~D notices are
electronic only from 2009; owners are matched by normalized name, which
succeeds for a minority favouring formally constituted entities; and the
final row conditions on funding realized after the filing being scored,
so it establishes that information distinguishing eventual listers was
present in the text years earlier, not that the language caused the
outcome.

Lead does not reverse anywhere after registration: negative at both use
proofs, at financing with or without registration imposed, and at listing
among the funded; the only row where it is not negative is SEC reporting
among registered owners, where it is indistinguishable from zero. Leading
filings fare worse at every stage where the market, rather than the
examiner, does the selecting.

Why would having been early cost? The two axes are not independent in the
obvious economic sense. What makes a filing leading is precisely that its
industry moved toward its language, and an industry that has moved toward
your language now describes its goods the way you describe yours. Being
early therefore consumes the quantity that pays: the mark that was
unusual in 2011 is ordinary by 2016 not because its owner changed
anything but because the category arrived. On this reading atypicality is
a position competitors erode, and lead measures how fast. The design
cannot demonstrate that mechanism; what can be said is that a negative
coefficient on lead at equal atypicality is what the mechanism would
produce if it were real.

A second reading follows from the first, and the design cannot test it
either. Privately, writing ahead of the industry's language is hard to
rationalize: no stage of the record rewards it, and each era's leading
vocabulary drew filers by the thousands anyway. Socially, the
exploration is not obviously waste. The corpus moves --- what makes a
filing leading is precisely that its class later wrote the way it did
--- so the language every industry eventually standardized on was first
filed by entrants who disproportionately did not survive to use it. If
finding the vocabulary an industry will adopt has value to the eventual
adopters, that value accrued to the followers while the cost fell on the
explorers.

\subsection{Limits}\label{sec:limits}

Three limits bound the reading. The vocabulary is positional rather than
causal: a leading filing was ahead of its industry's language, which
implies a trajectory without asserting that the filing caused it, and a
filing can sit where its class was heading because it was imitated,
because both responded to the same outside event, or by coincidence. What
drives the episodic pattern of Section~\ref{sec:waves} is not settled:
cohort and proof year move together, so post-crisis composition,
app-era speculative filings and calendar-time enforcement changes cannot
be separated on this design. And the registered description is a
conformed one: applicants who refile rewrite toward the class from both
tails, and an applicant who conformed before ever filing is unobserved,
so every estimate on registered marks is an estimate on
prosecution-filtered language, and the compression cannot be bounded
from below.

